\documentclass[a4paper]{article}

\usepackage[T1]{fontenc}
\usepackage[utf8]{inputenc}
\usepackage[ngerman]{babel}
\usepackage{amsmath,amssymb,amsthm}
\usepackage{tikz}
\usetikzlibrary{patterns}
\usepackage{float}
\usepackage{enumitem}
\usepackage[colorlinks=true,linkcolor=red,citecolor=red,urlcolor=red]{hyperref}

% --- Umgebungen (Stil des Kurzskripts: fettes Label, kursiver Text) ---------
\newtheoremstyle{kurzskript}%
  {\medskipamount}{\medskipamount}% Abstand oben / unten
  {\itshape}{}%                      Textschrift, Einzug
  {\bfseries}{}{ }%                  Labelschrift, Punkt, Abstand
  {\thmname{#1}\thmnumber{ #2}\thmnote{ \bfseries(#3)}}
\theoremstyle{kurzskript}
\newtheorem{satz}{Satz}[section]             % Nummerierung (1.1), (1.2), ...
\newtheorem{definition}[satz]{Definition}
\newtheorem{lemma}[satz]{Lemma}
\newtheorem{bemerkung}[satz]{Bemerkung}
\newtheorem{beispiel}[satz]{Beispiel}
\newtheorem{axiom}[satz]{Axiom}
\newtheorem{hilfssatz}[satz]{Hilfssatz}
\newtheorem*{axiomohne}{Axiom}
\newtheorem*{satzohne}{Satz}
\newtheorem*{bsp}{Beispiel}
\newtheorem*{bem}{Bemerkung}

\makeatletter
\renewenvironment{proof}[1][\proofname]{\par\pushQED{\qed}%
  \normalfont\topsep\medskipamount\relax
  \trivlist\item[\hskip\labelsep\bfseries #1\@addpunct{.}]\ignorespaces}%
  {\popQED\endtrivlist\@endpefalse}
\makeatother
\renewcommand{\proofname}{Beweis}
\renewcommand{\qedsymbol}{$\square$}

\newcommand{\N}{\mathbb{N}}
\newcommand{\Z}{\mathbb{Z}}
\newcommand{\Q}{\mathbb{Q}}
\newcommand{\R}{\mathbb{R}}
\newcommand{\Pm}{\mathbb{P}}
\newcommand{\Pot}{\mathcal{P}}

\title{Lineare Algebra 1}
\author{Mitschrift}
\date{Wintersemester 2026/27}

\begin{document}

\maketitle
\tableofcontents
\newpage

% ======================================================================
% Kapitel 1
% ======================================================================
\section*{Kapitel 1: Grundbegriffe -- Mengensprache und mathematische Argumentation}
\addcontentsline{toc}{part}{Kapitel 1: Grundbegriffe}

\paragraph{Untersuchungsgegenstände bzw.\ Werkzeuge der Linearen Algebra.}
Determinanten, Matrizen, Eigenwerte, Vektorräume, (lineare) Dimension, lineare Abbildungen, Multilinearformen usw.\ -- entwickelt vor allem im 18.\ und 19.~Jahrhundert.

\section{Mengen und Relationen}

Grundlage ist das \emph{Zermelo-Fraenkelsche Mengenaxiomensystem} (1907--1930).

\paragraph{Warum Axiome?}
\emph{Berrysches Paradoxon:} \glqq Welches ist die kleinste natürliche Zahl, die sich nicht mittels weniger als sechzehn deutschen Wörtern definieren lässt?\grqq{} -- Diese Beschreibung benötigt selbst nur 15 Wörter und definiert damit die Zahl, die sich angeblich nicht so definieren lässt.

Wir brauchen daher
\begin{itemize}
  \item eine formalisierte Sprache bzw.\ Schreibweise,
  \item mathematische Aussagen, die entweder wahr oder falsch sind -- nicht beides gleichzeitig.
\end{itemize}

\paragraph{Grundobjekte.}
Mengen $A, B, C, \dots, a, b, c, \dots$ und die Relation \glqq Element sein\grqq: $a \in A$ bzw.\ $a \notin A$. Man schreibt auch $A \ni a$, z.\,B.\ $\{1, 2, 3\} \ni 2$ und $\{1, 2, 3\} \not\ni 4$.

\begin{bemerkung}[\glqq bon mot\grqq{} -- Motivation nach Cantor]
Eine \emph{Menge} $A$ ist eine Zusammenfassung von wohlunterscheidbaren Objekten unserer Anschauung oder unseres Denkens, die \emph{Elemente} von $A$ genannt werden, zu einem Ganzen.
\end{bemerkung}

\begin{axiom}[Extensionalitätsaxiom]
Zwei Mengen $A$ und $B$ sind genau dann gleich, wenn sie dieselben Elemente haben. Kurz:
\[
  A = B \iff \forall x: \bigl(x \in A \iff x \in B\bigr).
\]
\end{axiom}

\begin{bem}
Für Mengen $M$ und $x$ ist genau eine der beiden Aussagen $x \in M$, $x \notin M$ wahr, die andere entsprechend falsch.
\end{bem}

Beim Bilden von \glqq neuartigen\grqq{} Mengen ist äußerste Vorsicht geboten. Eine häufig verwendete Methode liefert das

\begin{axiomohne}[Aussonderungsaxiom]
Zu jeder Menge $A$ und jeder Bedingung $P(x)$ an eine Menge $x$ gibt es die Menge $B$, deren Elemente genau die Elemente aus $A$ sind, die die Bedingung $P(x)$ erfüllen.

\emph{Schreibweise:} $B = \{\, x \in A \mid P(x) \,\}$.
\end{axiomohne}

\begin{bsp}
$\{\, n \in \N \mid \exists k \in \N : n = 2k \,\} = \{2, 4, 6, \dots\}$.
\end{bsp}

\begin{satzohne}[Russellsche Antinomie]
Zu jeder Menge $M$ gibt es eine Menge $x$ mit $x \notin M$. Insbesondere ist die Ansammlung aller Mengen so groß, dass sie \emph{keine} Menge bildet.
\end{satzohne}

\begin{proof}
Betrachte die \emph{Russell-Menge} (sie existiert nach dem Aussonderungsaxiom)
\[
  R(M) = \{\, x \in M \mid x \notin x \,\}.
\]
Angenommen, $R(M) \in M$.
\begin{itemize}
  \item Ist $R(M) \in R(M)$, so gilt nach Definition von $R(M)$ gerade $R(M) \notin R(M)$.
  \item Ist $R(M) \notin R(M)$, so erfüllt $R(M) \in M$ die Bedingung, also $R(M) \in R(M)$.
\end{itemize}
Beides ist ein Widerspruch. Also ist $R(M) \notin M$. Gäbe es eine Menge aller Mengen, so müsste sie $R(M)$ enthalten -- daher ist die Ansammlung aller Mengen keine Menge.
\end{proof}

\begin{beispiel}
\leavevmode
\begin{enumerate}[label=(\alph*)]
  \item \emph{Klassische Zahlbereiche:}
  \begin{align*}
    \N   &= \{1, 2, 3, \dots\}, \qquad \N_0 = \{0, 1, 2, 3, \dots\},\\
    \Z   &= \{\dots, -3, -2, -1, 0, 1, 2, 3, \dots\},\\
    \Q   &= \left\{ \tfrac{a}{b} \;\middle|\; a, b \in \Z,\ b \neq 0 \right\},
  \end{align*}
  $\R$ = reelle Zahlen, also alle Zahlen, die den Punkten auf der reellen Geraden entsprechen. \\
  (In dieser Vorlesung ist $0 \notin \N$.)

  \item Reihenfolge und Wiederholungen spielen keine Rolle (Extensionalität):
  \[
    \{1, 2, 3, 2\} = \{1, 2, 3\} = \left\{ \tfrac{4}{2},\ 5-2,\ (-1)^2 \right\}.
  \]

  \item Die \emph{leere Menge} $\emptyset = \{\}$ ist eindeutig bestimmt durch die Aussage
  \[
    \forall x: x \notin \emptyset.
  \]
  Mit ihr lassen sich die natürlichen Zahlen mengentheoretisch konstruieren:
  \[
    0 = \emptyset,\quad 1 = \{0\},\quad 2 = \{0, 1\},\quad \dots
  \]

  \item Die Menge der \emph{Primzahlen}:
  \[
    \Pm = \{\, p \in \N \mid p \text{ ist eine Primzahl} \,\} = \{2, 3, 5, 7, 11, 13, 17, \dots\}.
  \]
  Dabei gilt definitionsgemäß: $p \in \N$ ist eine Primzahl, falls $p \neq 1$ und gilt: Ist $p = ab$ das Produkt zweier natürlicher Zahlen $a, b$, so folgt $a = p,\ b = 1$ oder $a = 1,\ b = p$. Formal mit dem Aussonderungsaxiom:
  \[
    \Pm = \{\, p \in \N \mid p \neq 1 \ \land\ \forall a, b \in \N: \bigl(p = ab \implies (a = 1 \lor b = 1)\bigr) \,\}.
  \]
\end{enumerate}
\end{beispiel}

\begin{bemerkung}[Kernaufgaben der Mathematik]
\leavevmode
\begin{enumerate}
  \item Formuliere interessante Aussagen.
  \item Verifiziere diese mittels einer schlüssigen Beweisführung: \emph{direkte Beweise}, \emph{Widerspruchsbeweise}, \emph{Induktionsbeweise}.
\end{enumerate}
\end{bemerkung}

\begin{satz}[Beispiel für einen direkten Beweis]
Sei $n \in \N$. Dann gilt
\[
  \sum_{i=1}^{n} i = 1 + 2 + \dots + n = \frac{n(n+1)}{2}.
\]
\end{satz}

\begin{proof}
Mit den üblichen Rechenregeln für $\Q$ erhalten wir
\begin{align*}
  2 \cdot (1 + 2 + \dots + n)
  &= \bigl(1 + 2 + \dots + n\bigr) + \bigl(n + (n-1) + \dots + 1\bigr)\\
  &= \underbrace{(n+1) + (n+1) + \dots + (n+1)}_{n \text{ Summanden}} = n(n+1).
\end{align*}
Teilen durch $2$ liefert sofort die Behauptung.
\end{proof}

\begin{satz}[Euklid -- Beispiel für einen Widerspruchsbeweis]
Es gibt unendlich viele Primzahlen.
\end{satz}

\begin{proof}
Angenommen, die Menge der Primzahlen ist endlich, also $\Pm = \{p_1, p_2, \dots, p_r\}$. Betrachte
\[
  n = \prod_{i=1}^{r} p_i + 1.
\]
Es ist $n \geq 2$. Wir verwenden den folgenden Sachverhalt, den wir weiter unten begründen:

\smallskip
\noindent$(*)$ \emph{Ist $n \in \N$ mit $n \geq 2$, so besitzt $n$ einen Primteiler, d.\,h.\ es existiert eine Primzahl $p$ mit $p \mid n$, also $n = p \cdot m$ für ein geeignetes $m \in \N$.}
\smallskip

Sei also $n = p \cdot m$ mit $p \in \Pm$ und $m \in \N$. Ohne Einschränkung ist $p = p_1$. Dann erhalten wir
\[
  p_1 m = n = p_1 \prod_{i=2}^{r} p_i + 1,
  \qquad\text{also}\qquad
  p_1 \Bigl(m - \prod_{i=2}^{r} p_i\Bigr) = 1.
\]
Damit wäre $p_1 \mid 1$, aber wegen $p_1 \geq 2$ gilt $p_1 \nmid 1$. Widerspruch.

\smallskip
\emph{Zu $(*)$:} Angenommen, die Aussage ist falsch. Dann existiert ein kleinstes Gegenbeispiel $n \in \N$ mit $n \geq 2$, das keinen Primteiler besitzt. Offenbar ist dann $n \notin \Pm$ (sonst wäre $n$ selbst ein Primteiler), also $n = ab$ mit $a, b \in \N$ und $2 \leq a \leq b \leq n-1$. Nach Wahl von $n$ hat insbesondere $a$ einen Primteiler $p$. Wegen $p \mid a$ und $a \mid n$ gilt dann auch $p \mid n$ -- Widerspruch.
\end{proof}

\subsection*{Übung: Wahrheitstafeln}

Für Aussagen $\alpha, \beta$ definieren wir die Verknüpfungen durch ihre Wahrheitstafeln:
\begin{center}
\begin{tabular}{cc|c|c|c|c}
$\alpha$ & $\beta$ & $\alpha \land \beta$ & $\alpha \lor \beta$ & $\alpha \implies \beta$ & $\alpha \iff \beta$\\\hline
w & w & w & w & w & w\\
w & f & f & w & f & f\\
f & w & f & w & w & f\\
f & f & f & f & w & w
\end{tabular}
\end{center}

\begin{itemize}
  \item $\alpha \land \beta$ ist genau dann wahr, wenn $\alpha$ und $\beta$ beide wahr sind.
  \item $\alpha \lor \beta$ ist genau dann wahr, wenn mindestens eine der beiden Aussagen wahr ist.
  \item $\alpha \implies \beta$ ist genau dann falsch, wenn $\alpha$ wahr und $\beta$ falsch ist.
  \item $\alpha \iff \beta$ ist genau dann wahr, wenn $\alpha$ und $\beta$ denselben Wahrheitswert haben.
\end{itemize}

Zwei Aussagen heißen \emph{äquivalent}, falls sie die gleiche Wahrheitstafel besitzen.

\begin{bsp}
Die Aussagen $\neg\alpha \lor \neg\beta$ und $\alpha \implies \neg\beta$ sind äquivalent:
\begin{center}
\begin{tabular}{cc|cc|c|c}
$\alpha$ & $\beta$ & $\neg\alpha$ & $\neg\beta$ & $\neg\alpha \lor \neg\beta$ & $\alpha \implies \neg\beta$\\\hline
w & w & f & f & f & f\\
w & f & f & w & w & w\\
f & w & w & f & w & w\\
f & f & w & w & w & w
\end{tabular}
\end{center}
Beide Aussagen haben die gleiche Wahrheitstafel und sind damit äquivalent.
\end{bsp}

% ----------------------------------------------------------------------
\subsection{Elementare Mengenoperationen}

\begin{definition}[Elementare Mengenoperationen]
Seien $A$ und $B$ Mengen.
\begin{enumerate}[label=(\alph*)]
  \item $A$ heißt \emph{Teilmenge} von $B$, bez.\ $A \subseteq B$, falls gilt
  \[
    \forall x \in A: x \in B, \qquad\text{d.\,h.}\qquad \forall x: \bigl(x \in A \implies x \in B\bigr).
  \]
  $A$ heißt \emph{echte Teilmenge} von $B$, bez.\ $A \subsetneq B$, falls $A \subseteq B$, aber $A \neq B$. Auf Elementebene bedeutet das gerade:
  \[
    \forall x: \bigl(x \in A \implies x \in B\bigr) \quad\text{und}\quad \exists b \in B: b \notin A.
  \]
  \item Die \emph{Differenz} von $A$ und $B$ ist $A \setminus B = \{\, a \in A \mid a \notin B \,\}$.
  \item Die \emph{Potenzmenge} von $A$ ist $\Pot(A) = \{\, X \mid X \subseteq A \,\}$. Stets gilt $\emptyset \in \Pot(A)$ und $A \in \Pot(A)$.
  \item Die \emph{Vereinigung} von $A$ und $B$ ist $A \cup B = \{\, x \mid x \in A \lor x \in B \,\}$, \\
        der \emph{Schnitt} von $A$ und $B$ ist $A \cap B = \{\, x \mid x \in A \land x \in B \,\}$.
  \item Allgemeiner betrachten wir für eine beliebige Menge $\mathcal{A}$ von Mengen die \emph{große Vereinigung}
  \[
    \bigcup \mathcal{A} = \{\, x \mid \exists A \in \mathcal{A}: x \in A \,\}.
  \]
  In der Praxis ist häufig $\mathcal{A} = \{\, A_i \mid i \in I \,\}$ für eine \emph{Indexmenge} $I$, z.\,B.\ $I = \N$, und wir schreiben dann $\bigcup_{i \in I} A_i$.
  \item Ähnlich ist für $\mathcal{A} \neq \emptyset$ der \emph{große Schnitt}
  \[
    \bigcap \mathcal{A} = \bigcap_{A \in \mathcal{A}} A = \{\, x \mid \forall A \in \mathcal{A}: x \in A \,\}.
  \]
\end{enumerate}
\end{definition}

\begin{bem}
\leavevmode
\begin{itemize}
  \item $A = B$ ist gleichbedeutend mit $A \subseteq B \land B \subseteq A$.
  \item Es gilt $\Pm \subsetneq \N \subsetneq \Z \subsetneq \Q \subsetneq \R$.
\end{itemize}
\end{bem}

\begin{beispiel}
\leavevmode
\begin{enumerate}
  \item \emph{Cantor-Menge.} Setze
  \[
    C_1 = [0, 1], \qquad C_2 = \left[0, \tfrac13\right] \cup \left[\tfrac23, 1\right], \qquad
    C_n = \tfrac13 C_{n-1} \ \cup\ \left(\tfrac23 + \tfrac13 C_{n-1}\right).
  \]
  Man erhält eine absteigende Folge $C_1 \supsetneq C_2 \supsetneq C_3 \supsetneq \dots$ und definiert
  \[
    C = \bigcap_{n \in \N} C_n \subseteq [0, 1] \subseteq \R.
  \]

  \begin{figure}[H]
  \centering
  \begin{tikzpicture}[x=10cm,y=1cm]
    \foreach \n in {1,...,5} {
      \node[left] at (0,-0.7*\n) {$C_{\n}$};
    }
    \draw[line width=2.5pt] (0,-0.7) -- (1,-0.7);
    \foreach \a/\b in {0/1, 2/3} {
      \draw[line width=2.5pt] ({\a/3},-1.4) -- ({\b/3},-1.4);
    }
    \foreach \a in {0,2,6,8} {
      \draw[line width=2.5pt] ({\a/9},-2.1) -- ({(\a+1)/9},-2.1);
    }
    \foreach \a in {0,2,6,8,18,20,24,26} {
      \draw[line width=2.5pt] ({\a/27},-2.8) -- ({(\a+1)/27},-2.8);
    }
    \foreach \a in {0,2,6,8,18,20,24,26,54,56,60,62,72,74,78,80} {
      \draw[line width=2.5pt] ({\a/81},-3.5) -- ({(\a+1)/81},-3.5);
    }
    \node[below] at (0,-3.6) {$0$};
    \node[below] at (1,-3.6) {$1$};
  \end{tikzpicture}
  \caption{Die ersten Schritte der Konstruktion der Cantor-Menge}
  \end{figure}

  \emph{Frage:} Wie groß ist die Cantor-Menge $C$? \\
  Die Hausdorff-Dimension von $C$ ist $\dfrac{\log 2}{\log 3} \approx 0{,}63$.

  \item $\Pot(A)$ hat stets (deutlich) mehr Elemente als $A$ selbst [Satz von Cantor, siehe §3]. Für $A = \{1, 2, 3\}$ ist
  \[
    \Pot(A) = \bigl\{ \emptyset, \{1\}, \{2\}, \{3\}, \{1, 2\}, \{1, 3\}, \{2, 3\}, \{1, 2, 3\} \bigr\},
  \]
  also $|\Pot(A)| = 2^3 = 8$, während $|A| = 3$.
\end{enumerate}
\end{beispiel}

\begin{lemma}
Sei $M$ eine Menge und seien $A, B, C \subseteq M$. Dann gilt:
\begin{enumerate}[label=(\alph*)]
  \item $A \cup B = B \cup A$ \ und \ $A \cap B = B \cap A$ \hfill (Kommutativität)
  \item $A \cup (B \cup C) = (A \cup B) \cup C$ \ und \ $A \cap (B \cap C) = (A \cap B) \cap C$ \hfill (Assoziativität)
  \item $A \cap (B \cup C) = (A \cap B) \cup (A \cap C)$ \ und \ $A \cup (B \cap C) = (A \cup B) \cap (A \cup C)$ \hfill (Distributivität)
  \item $M \setminus (A \cup B) = (M \setminus A) \cap (M \setminus B)$ \ und \ $M \setminus (A \cap B) = (M \setminus A) \cup (M \setminus B)$ \hfill (de Morgansche Regeln)
\end{enumerate}
\end{lemma}

\begin{proof}
Die Teile (a), (b), (c) lassen sich \glqq leicht\grqq{} nachprüfen (z.\,B.\ mithilfe einer Wahrheitstafel). Wir beweisen den schwierigsten Teil (d).

Zunächst zeigen wir
\[
  (*) \qquad M \setminus (A \cup B) = (M \setminus A) \cap (M \setminus B).
\]
Dies lässt sich mit einem Venn-Diagramm veranschaulichen:

\begin{figure}[H]
\centering
\begin{tikzpicture}[scale=0.8]
  \def\rect{(-2.2,-1.6) rectangle (2.2,1.8)}
  \def\circA{(-0.55,0.25) circle (1)}
  \def\circB{(0.55,-0.15) circle (1)}
  % links: M \ A und M \ B
  \begin{scope}
    \fill[pattern=north east lines] \rect;
    \fill[white] \circA;
  \end{scope}
  \begin{scope}
    \begin{scope}[even odd rule]
      \clip \rect \circB;
      \fill[pattern=north west lines] \rect;
    \end{scope}
  \end{scope}
  \draw \rect; \draw \circA; \draw \circB;
  \node at (-0.85,0.45) {$A$};
  \node at (0.85,-0.35) {$B$};
  \node[above right] at (2.2,1.8) {$M$};
  \node[below] at (0,-1.7) {$(M \setminus A) \cap (M \setminus B)$: doppelt schraffiert};

  % rechts: M \ (A u B)
  \begin{scope}[xshift=6cm]
    \fill[gray!35] \rect;
    \fill[white] \circA;
    \fill[white] \circB;
    \draw \rect; \draw \circA; \draw \circB;
    \node at (-0.85,0.45) {$A$};
    \node at (0.85,-0.35) {$B$};
    \node[above right] at (2.2,1.8) {$M$};
    \node[below] at (0,-1.7) {$M \setminus (A \cup B)$};
  \end{scope}
\end{tikzpicture}
\caption{Veranschaulichung der de Morganschen Regel $(*)$}
\end{figure}

\glqq$\subseteq$\grqq: Sei $x \in M \setminus (A \cup B)$. Dann ist $x \in M$ und $x \notin A \cup B$, also $x \in M$ und $x \notin A$ und $x \notin B$. Also ist $x \in M \setminus A$ und $x \in M \setminus B$, d.\,h.\ $x \in (M \setminus A) \cap (M \setminus B)$.

\glqq$\supseteq$\grqq: Sei nun $x \in (M \setminus A) \cap (M \setminus B)$. Dann ist $x \in M$, $x \notin A$ und $x \notin B$. Also ist $x$ weder in $A$ noch in $B$, d.\,h.\ $x \notin A \cup B$, und damit $x \in M \setminus (A \cup B)$.

Damit ist $(*)$ bewiesen. Die zweite Aussage lässt sich direkt auf $(*)$ zurückführen. Dazu schreiben wir
\[
  \overline{C} = M \setminus C
\]
für das \emph{Komplement} von $C \subseteq M$ in $M$ und beobachten $\overline{\overline{C}} = C$. Wenden wir $(*)$ auf $\overline{A}$ und $\overline{B}$ anstelle von $A$ und $B$ an, so erhalten wir
\[
  \overline{\overline{A} \cup \overline{B}} = \overline{\overline{A}} \cap \overline{\overline{B}} = A \cap B.
\]
Bildet man auf beiden Seiten das Komplement, folgt
\[
  M \setminus (A \cap B) = \overline{A \cap B} = \overline{A} \cup \overline{B} = (M \setminus A) \cup (M \setminus B). \qedhere
\]
\end{proof}

% ----------------------------------------------------------------------
\subsection{Kartesisches Produkt}

\begin{definition}
\leavevmode
\begin{enumerate}[label=(\alph*)]
  \item Seien $A$ und $B$ Mengen. Das \emph{kartesische Produkt} von $A$ und $B$ ist
  \[
    A \times B = \{\, (a, b) \mid a \in A,\ b \in B \,\},
  \]
  die Menge aller \emph{geordneten Paare} $(a, b)$ mit $a \in A$ und $b \in B$. Diese werden mengentheoretisch so gebildet, dass für alle $x, y, \tilde{x}, \tilde{y}$ gilt:
  \[
    (x, y) = (\tilde{x}, \tilde{y}) \iff x = \tilde{x} \ \land\ y = \tilde{y}.
  \]
  Zum Beispiel kann man $(x, y) = \bigl\{ \{x\}, \{x, y\} \bigr\}$ setzen (Kuratowski).

  \item Für endlich viele Mengen $A_1, \dots, A_n$ ($n \in \N_0$) bilden wir entsprechend
  \[
    A_1 \times A_2 \times \dots \times A_n = \prod_{i=1}^{n} A_i = \{\, (a_1, a_2, \dots, a_n) \mid a_i \in A_i \text{ für } 1 \leq i \leq n \,\},
  \]
  wobei für $x_1, \dots, x_n, \tilde{x}_1, \dots, \tilde{x}_n$ stets gilt:
  \[
    (x_1, \dots, x_n) = (\tilde{x}_1, \dots, \tilde{x}_n) \iff x_1 = \tilde{x}_1 \ \land\ \dots\ \land\ x_n = \tilde{x}_n.
  \]

  \item Eine \emph{(binäre) Relation} auf einer Menge $A$ ist eine Teilmenge $R \subseteq A \times A$. Allgemeiner können Relationen auf einem Mengenpaar $A, B$ betrachtet werden: $R \subseteq A \times B$.
\end{enumerate}
\end{definition}

\begin{beispiel}
\leavevmode
\begin{enumerate}[label=(\alph*)]
  \item Die \emph{euklidische Ebene} $\mathbb{E}_2 = \R^2 = \R \times \R$ ist die Koordinatenebene:
  \begin{figure}[H]
  \centering
  \begin{tikzpicture}[scale=0.8,>=stealth]
    \draw[->] (-0.5,0) -- (5,0) node[right] {$x$};
    \draw[->] (0,-0.5) -- (0,3.5) node[above] {$y$};
    \fill (3,2) circle (2pt) node[above right] {$(a, b)$};
    \draw[dashed] (3,0) node[below] {$a$} -- (3,2) -- (0,2) node[left] {$b$};
  \end{tikzpicture}
  \caption{$\R \times \R$}
  \end{figure}

  \item Die Menge $\{\, (a, a) \mid a \in \R \,\} \subseteq \R^2$ kann geometrisch als \emph{Gerade} (die Winkelhalbierende) interpretiert werden -- oder auch als Relation auf $\R$, nämlich als die \emph{Gleichheitsrelation} $=_\R$ für reelle Zahlen.

  \item Die Menge $\{\, (a, a + b^2) \mid a, b \in \R \,\} = \{\, (x, y) \in \R^2 \mid x \leq y \,\}$ ist die \emph{Kleinergleichrelation} $\leq_\R$ auf $\R$. Geometrisch ist das die Halbebene oberhalb der Winkelhalbierenden (einschließlich dieser).
\end{enumerate}

\begin{figure}[H]
\centering
\begin{tikzpicture}[scale=0.55,>=stealth]
  % links: Gleichheitsrelation
  \draw[gray!30,very thin] (-3,-3) grid (3,3);
  \draw[->] (-3.3,0) -- (3.5,0) node[right] {$x$};
  \draw[->] (0,-3.3) -- (0,3.5) node[above] {$y$};
  \draw[thick,green!50!black] (-3,-3) -- (3,3);
  \node[below] at (0,-3.4) {$=_\R$};

  % rechts: Kleinergleichrelation
  \begin{scope}[xshift=9cm]
    \fill[blue!20] (-3,-3) -- (3,3) -- (-3,3) -- cycle;
    \draw[gray!30,very thin] (-3,-3) grid (3,3);
    \draw[->] (-3.3,0) -- (3.5,0) node[right] {$x$};
    \draw[->] (0,-3.3) -- (0,3.5) node[above] {$y$};
    \draw[thick,blue] (-3,-3) -- (3,3);
    \node[below] at (0,-3.4) {$\leq_\R$};
  \end{scope}
\end{tikzpicture}
\caption{Gleichheitsrelation und Kleinergleichrelation als Teilmengen von $\R^2$}
\end{figure}
\end{beispiel}

% ======================================================================
\section{Abbildungen und Bijektionen}

\begin{bemerkung}[Motivation]
Eine Abbildung von $A$ nach $B$ ordnet jedem $a \in A$ ein Element $b_a \in B$ zu: $A \to B$, $a \mapsto b_a$.
\begin{enumerate}[label=(\arabic*)]
  \item $x \mapsto x^2$ für $A = B = \R$. Diese Zuordnung lässt sich mit ihrem Graphen $\{\, (x, x^2) \mid x \in \R \,\} \subseteq \R \times \R$ identifizieren:
  \begin{figure}[H]
  \centering
  \begin{tikzpicture}[scale=0.5,>=stealth]
    \draw[gray!30,very thin] (-4,-1) grid (4,7);
    \draw[->] (-4.3,0) -- (4.5,0) node[right] {$x$};
    \draw[->] (0,-1.3) -- (0,7.5) node[above] {$y$};
    \foreach \t in {-2,2} \node[below] at (\t,0) {\small $\t$};
    \foreach \t in {2,4,6} \node[left] at (0,\t) {\small $\t$};
    \draw[thick,domain=-2.6:2.6,samples=60] plot (\x,{\x*\x});
  \end{tikzpicture}
  \caption{Graph von $x \mapsto x^2$}
  \end{figure}

  \item Mit $\pi = 3{,}141592653589\ldots$ betrachte
  \[
    \N \to \{0, 1, 2, \dots, 9\}, \qquad n \mapsto n\text{-te Nachkommastelle von } \pi,
  \]
  also $1 \mapsto 1$, $2 \mapsto 4$, \dots:
  \begin{center}
  \begin{tabular}{c|ccccc}
  $n$ & 1 & 2 & 3 & 4 & 5 \\\hline
  Bild & 1 & 4 & 1 & 5 & 9
  \end{tabular}
  \end{center}
\end{enumerate}
\end{bemerkung}

\begin{definition}
Seien $A, B$ Mengen. Eine \emph{Abbildung} $\alpha$ von $A$ nach $B$, i.\,Z.\ $\alpha\colon A \to B$, ist eine Relation $\alpha \subseteq A \times B$ mit der (wesentlichen) Eigenschaft:
\[
  \text{Für jedes } a \in A \text{ gibt es genau ein } b \in B \text{ mit } (a, b) \in \alpha.
\]
Für eine Abbildung $\alpha\colon A \to B$ heißt $A$ der \emph{Definitionsbereich} und $B$ \emph{ein Wertebereich} von $\alpha$. Ist $(a, b) \in \alpha$, so schreiben wir prägnanter
\[
  a\alpha = b, \qquad a.\alpha = b, \qquad (a)\alpha = b \qquad\text{oder}\qquad a^\alpha = b
\]
und nennen $a\alpha$ das \emph{Bildelement} von $a$ unter $\alpha$. Schließlich heißt
\[
  \operatorname{Bild}(\alpha) = A\alpha = \{\, a\alpha \mid a \in A \,\} \subseteq B
\]
das (gesamte) \emph{Bild} von $\alpha$.
\end{definition}

\begin{bem}
\leavevmode
\begin{itemize}
  \item Achtung: In der Analysis würden wir fast immer $\alpha(a) = b$ schreiben. In dieser Vorlesung stehen Abbildungen \emph{rechts} vom Argument.
  \item Für $\R \to \R$, $x \mapsto x^2$ ist das Bild gerade $\R_{\geq 0}$.
  \item Da eine Abbildung als Relation definiert ist, ist der Wertebereich nicht eindeutig festgelegt: Für uns sind die Abbildungen $\sin\colon \R \to \R$ und $\sin\colon \R \to [-1, 1]$ gleich.
  \item Die \emph{leere Abbildung} ist $\emptyset\colon \emptyset \to \emptyset$ (allgemeiner $\emptyset\colon \emptyset \to B$ für jede Menge $B$).
\end{itemize}
\end{bem}

\begin{definition}
Seien $A, B$ Mengen.
\begin{enumerate}[label=(\alph*)]
  \item $\operatorname{Abb}(A, B) = \{\, \alpha \mid \alpha\colon A \to B \,\}$ ist die Menge aller Abbildungen von $A$ nach $B$.
  \item Eine Abbildung $\alpha\colon A \to B$ heißt \emph{injektiv}, falls aus $a\alpha = \tilde{a}\alpha$ stets $a = \tilde{a}$ folgt.
  \item Die Abbildung $\alpha$ bildet $A$ \emph{surjektiv} auf $B$ ab, falls $\operatorname{Bild}(\alpha) = B$.
\end{enumerate}
Schließlich heißt $\alpha\colon A \to B$ eine \emph{Bijektion} von $A$ auf $B$ (kurz: \emph{bijektiv}), falls $\alpha$ injektiv ist und $A$ surjektiv auf $B$ abbildet.
\end{definition}

\begin{bem}
Mit zunehmender Erfahrung schreiben wir (auch) kurz: \glqq$\alpha\colon A \to B$ ist surjektiv bzw.\ bijektiv.\grqq
\end{bem}

\begin{beispiel}
\leavevmode
\begin{enumerate}[label=(\alph*)]
  \item Die Abbildung $f\colon \R \to \R$, $x \mapsto x^2$ ist weder injektiv (z.\,B.\ $(-2)f = 4 = 2f$) noch surjektiv auf $\R$ (z.\,B.\ $-1 \notin \operatorname{Bild}(f)$). \\
  Ihre Einschränkung auf $\R_{\geq 0}$, d.\,h.\ die Abbildung $g\colon \R_{\geq 0} \to \R_{\geq 0}$, $x \mapsto x^2$, ist eine Bijektion von $\R_{\geq 0}$ auf sich.
  \begin{figure}[H]
  \centering
  \begin{tikzpicture}[scale=0.5,>=stealth]
    \draw[gray!30,very thin] (-4,-1) grid (4,7);
    \draw[->] (-4.3,0) -- (4.5,0) node[right] {$x$};
    \draw[->] (0,-1.3) -- (0,7.5) node[above] {$y$};
    \draw[thick,gray,domain=-2.6:0,samples=40] plot (\x,{\x*\x});
    \draw[ultra thick,blue,domain=0:2.6,samples=40] plot (\x,{\x*\x}) node[right] {$g$};
    \draw[dashed] (-2,0) node[below] {\small $-2$} -- (-2,4) -- (2,4) -- (2,0) node[below] {\small $2$};
    \node[left] at (0,4) {\small $4$};
  \end{tikzpicture}
  \caption{$f$ (grau und blau) und ihre Einschränkung $g$ auf $\R_{\geq 0}$ (blau)}
  \end{figure}

  \item Seien $a, b \in \R$ und
  \[
    \alpha = \alpha_{a, b}\colon \R \to \R, \qquad x \mapsto ax + b.
  \]
  \emph{Behauptung:} Für $a \neq 0$ ist $\alpha_{a, b}$ eine Bijektion von $\R$ auf sich.
  \begin{proof}
  Sei $a \neq 0$.

  \emph{Injektivität:} Seien $x, y \in \R$ mit $x\alpha = y\alpha$. Dann gilt $ax + b = ay + b$, also $ax = ay$, und wegen $a \neq 0$ folgt $x = y$.

  \emph{Surjektivität:} Sei $y \in \R$. Gesucht ist ein $x \in \R$ mit $x\alpha = y$. Setze $x = a^{-1}(y - b)$. Dann ist tatsächlich
  \[
    x\alpha = a\bigl(a^{-1}(y - b)\bigr) + b = y. \qedhere
  \]
  \end{proof}

  \item Für beliebige Mengen $A$ gilt
  \[
    \operatorname{Abb}(\emptyset, A) = \{\emptyset\}, \qquad \operatorname{Abb}(A, \emptyset) = \emptyset \ \text{ für } A \neq \emptyset.
  \]
  Die Menge $\operatorname{Abb}(\{1, 2, 3\}, \{0, 1\})$ besteht aus genau $8 = 2^3$ Elementen. In solchen Situationen verwenden wir häufig \emph{Abbildungstafeln}; z.\,B.\ bezeichnet
  \[
    \begin{pmatrix} 1 & 2 & 3 \\ 0 & 1 & 1 \end{pmatrix}
  \]
  die Abbildung $\{1, 2, 3\} \to \{0, 1\}$, die durch $1 \mapsto 0$, $2 \mapsto 1$ und $3 \mapsto 1$ gegeben ist. Streng genommen ist diese Abbildung die Menge
  \[
    \{(1, 0), (2, 1), (3, 1)\}.
  \]

  \item Sei $A$ eine Menge. Für $T \subseteq A$ bezeichne
  \[
    \iota_T\colon A \to \{0, 1\}, \qquad a \mapsto
    \begin{cases}
      0 & \text{falls } a \notin T,\\
      1 & \text{falls } a \in T
    \end{cases}
  \]
  die zugehörige \emph{Indikatorfunktion}. Dann liefert
  \[
    \iota\colon \Pot(A) \to \operatorname{Abb}(A, \{0, 1\}), \qquad T \mapsto \iota_T
  \]
  eine Bijektion. Mithilfe von $\iota$ lassen sich die einfachen \glqq Rechenregeln\grqq{} für $\cup$, $\cap$, $\setminus$ übersetzen (und damit verständlich machen) in entsprechende Rechenregeln für $\operatorname{Abb}(A, \{0, 1\})$, wobei mit den jeweiligen Bildern modulo~$2$ gerechnet wird. Etwa gilt für alle $a \in A$:
  \[
    a\iota_{S \cap T} = (a\iota_S)(a\iota_T), \qquad a\iota_{A \setminus T} = 1 + a\iota_T \pmod 2.
  \]
\end{enumerate}
\end{beispiel}

\begin{satz}[und Definition]
\leavevmode
\begin{enumerate}[label=(\alph*)]
  \item Seien $A, B, C, D$ Mengen und $\alpha\colon A \to B$, $\beta\colon B \to C$, $\gamma\colon C \to D$ Abbildungen. Das \emph{Kompositum} (oder: die \emph{Hintereinanderausführung}) von $\alpha$ und $\beta$ ist die Abbildung
  \[
    \alpha\beta\colon A \to C, \qquad a \mapsto (a\alpha)\beta.
  \]
  \begin{center}
  \begin{tikzpicture}[>=stealth,node distance=2.4cm]
    \node (A) {$A$};
    \node[right of=A] (B) {$B$};
    \node[right of=B] (C) {$C$};
    \node[right of=C] (D) {$D$};
    \draw[->] (A) -- node[above] {$\alpha$} (B);
    \draw[->] (B) -- node[above] {$\beta$} (C);
    \draw[->] (C) -- node[above] {$\gamma$} (D);
    \draw[->] (A) to[bend right=35] node[below] {$\alpha\beta$} (C);
    \draw[->] (B) to[bend left=35] node[above] {$\beta\gamma$} (D);
  \end{tikzpicture}
  \end{center}
  Es gilt das \emph{Assoziativgesetz}
  \[
    (\alpha\beta)\gamma = \alpha(\beta\gamma).
  \]
  \item Die \emph{identische Abbildung} auf einer Menge $A$ ist $\operatorname{id}_A\colon A \to A$, $a \mapsto a$. Für jede Abbildung $\alpha\colon A \to B$ gilt
  \[
    \operatorname{id}_A \alpha = \alpha = \alpha \operatorname{id}_B.
  \]
  \item Sei $\alpha\colon A \to B$ injektiv. Dann ist
  \[
    \alpha^{-1} = \{\, (b, a) \mid (a, b) \in \alpha \,\} \subseteq \operatorname{Bild}(\alpha) \times A
  \]
  eine Abbildung $\alpha^{-1}\colon \operatorname{Bild}(\alpha) \to A$ mit
  \[
    \alpha\alpha^{-1} = \operatorname{id}_A \qquad\text{und}\qquad \alpha^{-1}\alpha = \operatorname{id}_{\operatorname{Bild}(\alpha)}.
  \]
  Die Abbildung $\alpha^{-1}$ heißt entsprechend die \emph{Umkehrabbildung} zu $\alpha$.
\end{enumerate}
\end{satz}

\begin{proof}
(a) Beide Seiten sind Abbildungen von $A$ nach $D$, und für beliebiges $a \in A$ gilt
\[
  a\bigl((\alpha\beta)\gamma\bigr) = \bigl(a(\alpha\beta)\bigr)\gamma = \bigl((a\alpha)\beta\bigr)\gamma = (a\alpha)(\beta\gamma) = a\bigl(\alpha(\beta\gamma)\bigr).
\]
(b) Für $a \in A$ ist $a(\operatorname{id}_A\alpha) = (a \operatorname{id}_A)\alpha = a\alpha$ und $a(\alpha\operatorname{id}_B) = (a\alpha)\operatorname{id}_B = a\alpha$.

(c) Um einzusehen, dass $\alpha^{-1}$ eine Abbildung ist, sei $b \in \operatorname{Bild}(\alpha)$. Dann gibt es ein $a \in A$ mit $a\alpha = b$, also $(b, a) \in \alpha^{-1}$. Seien nun $(b, a), (b, \tilde{a}) \in \alpha^{-1}$; zu zeigen ist $a = \tilde{a}$. Da $\alpha$ injektiv ist, folgt aus $a\alpha = b = \tilde{a}\alpha$ tatsächlich $a = \tilde{a}$. Damit ist $\alpha^{-1}$ als Abbildung mit Definitionsbereich $\operatorname{Bild}(\alpha) \subseteq B$ und Wertebereich (sogar Bild) $A$ erkannt.

Die weiteren Behauptungen rechnen wir sofort nach: Für $a \in A$ ist $(a, a\alpha) \in \alpha$, also $(a\alpha)\alpha^{-1} = a$, d.\,h.\ $\alpha\alpha^{-1} = \operatorname{id}_A$. Für $b \in \operatorname{Bild}(\alpha)$ ist $b\alpha^{-1}$ gerade das (eindeutige) $a$ mit $a\alpha = b$, also $(b\alpha^{-1})\alpha = b$, d.\,h.\ $\alpha^{-1}\alpha = \operatorname{id}_{\operatorname{Bild}(\alpha)}$. Insbesondere ist auch $\operatorname{Bild}(\alpha^{-1}) = A$.
\end{proof}

\begin{bem}
Für die Hintereinanderausführung $\alpha\beta$ haben wir implizit vorausgesetzt, dass $\operatorname{Bild}(\alpha)$ im Definitionsbereich von $\beta$ liegt. Verzichtet man auf diese Bedingung, ist also $\beta\colon B' \to C$ mit einer beliebigen Menge $B'$, so ergibt sich trotzdem eine Abbildung
\[
  \alpha\beta\colon A_0 \to C \qquad\text{mit}\qquad A_0 = \{\, a \in A \mid a\alpha \in B' \,\}
\]
mit entsprechend eingeschränktem Definitionsbereich.
\end{bem}

\begin{lemma}
Für eine Abbildung $\alpha\colon A \to B$ gelten:
\begin{enumerate}[label=(\alph*)]
  \item $\alpha$ ist injektiv $\iff$ $\exists B_0 \subseteq B\ \exists \beta \in \operatorname{Abb}(B_0, A)$: $\alpha\beta = \operatorname{id}_A$.
  \item $\alpha$ ist surjektiv (auf $B$) $\iff$ $\exists \gamma \in \operatorname{Abb}(B, A)$: $\gamma\alpha = \operatorname{id}_B$.
\end{enumerate}
\end{lemma}

\begin{bem}
Ist $A \neq \emptyset$, so bleibt (a) richtig, wenn wir starr $B_0 = B$ setzen. Dann reicht $\beta = \alpha^{-1}$ wie im Beweis unten nicht immer aus, da $\alpha^{-1}$ nur auf $\operatorname{Bild}(\alpha)$ definiert ist; man setzt $\alpha^{-1}$ dann beliebig auf $B$ fort, etwa durch $b \mapsto a_0$ für alle $b \in B \setminus \operatorname{Bild}(\alpha)$ mit einem festen $a_0 \in A$. (Für $A = \emptyset \neq B$ geht das nicht, denn dann ist $\operatorname{Abb}(B, \emptyset) = \emptyset$.)
\end{bem}

\begin{proof}
(a) \glqq$\Rightarrow$\grqq: Ist $\alpha$ injektiv, so wählen wir $B_0 = \operatorname{Bild}(\alpha)$ und $\beta = \alpha^{-1}$ (Satz~2.5\,(c)).

\glqq$\Leftarrow$\grqq: Seien nun $B_0 \subseteq B$ und $\beta\colon B_0 \to A$ mit $\alpha\beta = \operatorname{id}_A$ gegeben. Zu zeigen ist, dass $\alpha$ injektiv ist. Seien $a, \tilde{a} \in A$ mit $a\alpha = \tilde{a}\alpha$. Dann folgt tatsächlich
\[
  a = a\operatorname{id}_A = a(\alpha\beta) = (a\alpha)\beta = (\tilde{a}\alpha)\beta = \tilde{a}(\alpha\beta) = \tilde{a}\operatorname{id}_A = \tilde{a}.
\]

(b) \glqq$\Rightarrow$\grqq: Sei zunächst $\alpha$ surjektiv auf $B$. Zu jedem $b \in B$ finden wir ein $a_b \in A$ mit $a_b\alpha = b$. Wir definieren $\gamma\colon B \to A$ mittels $b\gamma = a_b$.\,$(*)$ Für beliebiges $b \in B$ gilt dann
\[
  b(\gamma\alpha) = (b\gamma)\alpha = a_b\alpha = b,
\]
also ist $\gamma\alpha = \operatorname{id}_B$.

\glqq$\Leftarrow$\grqq: Sei umgekehrt $\gamma\colon B \to A$ mit $\gamma\alpha = \operatorname{id}_B$ gegeben. Zu zeigen ist, dass $\alpha$ surjektiv auf $B$ ist. Sei $b \in B$ gegeben und setze $a = b\gamma \in A$. Dann gilt tatsächlich
\[
  b = b\operatorname{id}_B = b(\gamma\alpha) = (b\gamma)\alpha = a\alpha \in \operatorname{Bild}(\alpha). \qedhere
\]
\end{proof}

\begin{bem}
Für $(*)$, also die gleichzeitige Auswahl eines Urbilds $a_b$ für \emph{jedes} $b \in B$, wird das \emph{Auswahlaxiom} verwendet.
\end{bem}

\begin{hilfssatz}
Eine Abbildung $\alpha\colon A \to B$ ist genau dann eine Bijektion von $A$ auf $B$, wenn Abbildungen $\beta, \gamma\colon B \to A$ existieren mit
\[
  \alpha\beta = \operatorname{id}_A \qquad\text{und}\qquad \gamma\alpha = \operatorname{id}_B.
\]
In diesem Fall gilt bereits $B = \operatorname{Bild}(\alpha)$ und $\beta = \gamma = \alpha^{-1}$.
\end{hilfssatz}

\begin{proof}
\glqq$\Rightarrow$\grqq: Sei $\alpha\colon A \to B$ bijektiv. Insbesondere bildet $\alpha$ surjektiv auf $B$ ab, d.\,h.\ $\operatorname{Bild}(\alpha) = B$, und somit liefern Satz~2.5\,(c) und Lemma~2.6 die Abbildung $\beta = \alpha^{-1}\colon B \to A$ sowie $\gamma$ mit $\alpha\beta = \operatorname{id}_A$ und $\gamma\alpha = \operatorname{id}_B$.

\glqq$\Leftarrow$\grqq: Sind umgekehrt $\beta, \gamma\colon B \to A$ gegeben mit $\alpha\beta = \operatorname{id}_A$ und $\gamma\alpha = \operatorname{id}_B$, so liefert Lemma~2.6 direkt: $\alpha$ ist injektiv und surjektiv auf $B$.

\emph{Zusatz:} Wegen $\gamma\alpha = \operatorname{id}_B$ ist $\operatorname{Bild}(\alpha) = B$ (Lemma~2.6\,(b)). Weiter ist
\[
  \beta = \operatorname{id}_B\beta = (\gamma\alpha)\beta = \gamma(\alpha\beta) = \gamma\operatorname{id}_A = \gamma.
\]
Die Bedingungen $\alpha\beta = \operatorname{id}_A$ und $\beta\alpha = \operatorname{id}_B$ liefern dann $\gamma = \beta = \alpha^{-1}$: Für $b \in B$ gibt es nämlich $a \in A$ mit $a\alpha = b$, und es folgt $b\beta = (a\alpha)\beta = a = b\alpha^{-1}$.
\end{proof}

% ======================================================================
\section{Abzählen und Mächtigkeit von Mengen}

\begin{definition}
\leavevmode
\begin{enumerate}[label=(\alph*)]
  \item Zwei Mengen $A, B$ heißen \emph{gleichmächtig} (auch: von gleicher \emph{Mächtigkeit} bzw.\ \emph{Kardinalität}), falls es eine Bijektion $f\colon A \to B$ gibt. Wir schreiben dann $A \approx B$.

  \emph{Merke:} Offenbar gelten für Mengen $A, B, C$:
  \begin{enumerate}[label=(\roman*)]
    \item $A \approx A$;
    \item $A \approx B \iff B \approx A$;
    \item aus $A \approx B$ und $B \approx C$ folgt $A \approx C$.
  \end{enumerate}
  (Verwende $\operatorname{id}_A$, die Umkehrabbildung $f^{-1}$ bzw.\ das Kompositum zweier Bijektionen.)

  \item Eine Menge $A$ heißt \emph{endlich}, falls $A \approx \{1, 2, \dots, n\}$ für ein geeignetes $n \in \N_0$ gilt (für $n = 0$ ist $\{1, \dots, n\} = \emptyset$). Ist in diesem Fall $f\colon \{1, 2, \dots, n\} \to A$ eine Bijektion, so gilt offenbar
  \[
    A = \{a_1, a_2, \dots, a_n\},
  \]
  wobei die Elemente $a_i = if$, $1 \leq i \leq n$, paarweise verschieden sind. Ist $A$ nicht endlich, so heißt $A$ \emph{unendlich}; z.\,B.\ ist $\N$ unendlich.

  Für $m, n \in \N_0$ gilt: $\{1, 2, \dots, m\} \approx \{1, 2, \dots, n\} \iff m = n$.
\end{enumerate}
Dementsprechend definieren wir die \emph{Mächtigkeit} einer Menge $A$ als
\[
  \#A = |A| =
  \begin{cases}
    n & \text{falls } A \approx \{1, 2, \dots, n\} \text{ für ein } n \in \N_0,\\
    \infty & \text{andernfalls.}
  \end{cases}
\]
Zur Illustration unterscheiden wir bestenfalls noch zwischen
\begin{itemize}
  \item $A \approx \N$: $A$ ist \emph{abzählbar unendlich}, und
  \item $A$ ist unendlich, aber nicht abzählbar unendlich: $A$ ist \emph{überabzählbar}.
\end{itemize}
\end{definition}

\begin{beispiel}
\leavevmode
\begin{enumerate}[label=(\alph*)]
  \item Es gilt
  \[
    \N \subsetneq \N_0 \subsetneq \Z \subsetneq \Q \subsetneq \R,
  \]
  aber $\N \approx \N_0 \approx \Z \approx \Q$ und $\N \not\approx \R$. Zum Beispiel ist
  \[
    \N_0 \to \N, \qquad n \mapsto n + 1
  \]
  eine Bijektion (mit Umkehrabbildung $\N \to \N_0$, $n \mapsto n - 1$).

  \item $\N \approx \Z \approx \Q$. Eine Bijektion $\N \to \Z$ ist gegeben durch die Abbildungstafel
  \[
    \begin{pmatrix} 1 & 2 & 3 & 4 & 5 & 6 & \dots \\ 0 & 1 & -1 & 2 & -2 & 3 & \dots \end{pmatrix}
  \]
  (\glqq Reißverschlusstrick\grqq). Für $\Q$ kann mit dem gleichen Reißverschlusstrick die Behauptung $\N \approx \Q$ auf $\N \approx \Q^+$ zurückgeführt werden. Jede positive rationale Zahl können wir als gekürzten Bruch schreiben:
  \begin{figure}[H]
  \centering
  \begin{tikzpicture}[x=1.3cm,y=0.9cm]
    % Diagonalen a + b = i + 1
    \foreach \s in {2,...,6}
      \draw[blue!40,thick] (\s-1,-1) -- (1,-\s+1);
    \foreach \a in {1,...,5}
      \foreach \b in {1,...,5} {
        \pgfmathtruncatemacro{\g}{gcd(\a,\b)}
        \ifnum\g=1
          \node[fill=white,inner sep=1.5pt] at (\a,-\b) {$\frac{\a}{\b}$};
        \else
          \node[fill=white,inner sep=1.5pt,text=gray!60] at (\a,-\b) {$\frac{\a}{\b}$};
        \fi
      }
    \foreach \a in {1,...,5} \node at (\a,-5.8) {$\vdots$};
    \foreach \b in {1,...,5} \node at (5.8,-\b) {$\cdots$};
    \foreach \i in {1,...,5} \node[blue,font=\small] at (\i+0.45,-0.45) {$A_{\i}$};
  \end{tikzpicture}
  \caption{Positive Brüche $\frac{a}{b}$; auf der $i$-ten Diagonalen gilt $a + b = i + 1$. Nicht gekürzte Brüche sind grau.}
  \end{figure}
  Somit ist
  \[
    \Q^+ = \Bigl\{\, \frac{a}{b} \Bigm| a, b \in \N \text{ teilerfremd} \,\Bigr\} = \bigcup_{i \in \N} A_i
  \]
  mit paarweise disjunkten Teilmengen
  \[
    A_i = \Bigl\{\, \frac{a}{b} \Bigm| a, b \in \N \text{ teilerfremd mit } a + b = i + 1 \,\Bigr\},
  \]
  also etwa $A_1 = \{\frac{1}{1}\}$, $A_2 = \{\frac{1}{2}, \frac{2}{1}\}$, $A_3 = \{\frac{1}{3}, \frac{3}{1}\}$, $A_4 = \{\frac{1}{4}, \frac{2}{3}, \frac{3}{2}, \frac{4}{1}\}$. Jedes $A_i$ ist endlich mit höchstens $i$ Elementen.

  \item $\N \not\approx \R$. Zunächst liefert die Abbildung
  \[
    (0, 1)_\R \to \R, \qquad t \mapsto
    \begin{cases}
      \dfrac{1 - |2t - 1|}{2t - 1} & \text{falls } t \neq \tfrac12,\\[1ex]
      0 & \text{falls } t = \tfrac12
    \end{cases}
  \]
  eine Bijektion, also $(0, 1)_\R \approx \R$. Für $t < \frac12$ ist der Ausdruck gleich $\frac{2t}{2t - 1}$ und durchläuft $(-\infty, 0)$, für $t > \frac12$ ist er gleich $\frac{2 - 2t}{2t - 1}$ und durchläuft $(0, \infty)$.
  \begin{figure}[H]
  \centering
  \begin{tikzpicture}[x=6cm,y=0.45cm,>=stealth]
    \draw[gray!30,very thin,xstep=0.25,ystep=1] (0,-6) grid (1,6);
    \draw[->] (-0.1,0) -- (1.12,0) node[right] {$t$};
    \draw[->] (0,-6.3) -- (0,6.6);
    \foreach \y in {-4,-2,2,4} \node[left,font=\small] at (0,\y) {$\y$};
    \draw[dashed,gray] (0.5,-6) -- (0.5,6);
    \node[below,font=\small] at (0.5,-6) {$\tfrac12$};
    \node[below left,font=\small] at (0,0) {$0$};
    \node[below right,font=\small] at (1,0) {$1$};
    \draw[thick,green!50!black,domain=0.005:0.429,samples=80] plot (\x,{2*\x/(2*\x-1)});
    \draw[thick,green!50!black,domain=0.571:0.995,samples=80] plot (\x,{(2-2*\x)/(2*\x-1)});
    \fill[green!50!black] (0.5,0) circle (1.6pt);
    \draw[green!50!black,fill=white] (0,0) circle (1.6pt);
    \draw[green!50!black,fill=white] (1,0) circle (1.6pt);
  \end{tikzpicture}
  \caption{Eine Bijektion $(0, 1)_\R \to \R$}
  \end{figure}
  Es bleibt $(0, 1)_\R \not\approx \N$ zu zeigen; dazu dient das \emph{Cantorsche Diagonalargument}. Wir verwenden: Jede reelle Zahl $x \in (0, 1)_\R$ lässt sich \emph{eindeutig} als Dezimalbruch schreiben, wenn keine Perioden der Form $\dots 999\dots$ verwendet werden (z.\,B.\ $0{,}32999\ldots = 0{,}33$).
  \begin{proof}
  Angenommen, $(0, 1)_\R \approx \N$. Dann lassen sich alle Elemente von $(0, 1)_\R$ in einer Liste aufzählen:
  \begin{align*}
    x_1 &= 0{,}x_{11}x_{12}x_{13}\dots\\
    x_2 &= 0{,}x_{21}x_{22}x_{23}\dots\\
    x_3 &= 0{,}x_{31}x_{32}x_{33}\dots\\
    x_4 &= 0{,}x_{41}x_{42}x_{43}\dots\\
        &\ \ \vdots
  \end{align*}
  mit Ziffern $x_{ij} \in \{0, 1, \dots, 9\}$ (ohne $9$er-Perioden).
  \begin{figure}[H]
  \centering
  \begin{tikzpicture}[x=1.1cm,y=0.65cm]
    \foreach \i in {1,...,5} {
      \node[left] at (0.35,-\i) {$x_{\i} = 0{,}$};
      \foreach \j in {1,...,5} {
        \ifnum\i=\j
          \node[draw=red,thick,circle,inner sep=0.5pt,text=red] at (\j,-\i) {$x_{\i\j}$};
        \else
          \node at (\j,-\i) {$x_{\i\j}$};
        \fi
      }
      \node at (5.9,-\i) {$\cdots$};
    }
    \foreach \j in {1,...,5} \node at (\j,-5.8) {$\vdots$};
    \node at (5.9,-5.8) {$\ddots$};
    \draw[red!50,thick,dashed] (0.6,-0.6) -- (5.6,-5.6);
    \node[left,red] at (0.35,-7) {$y = 0{,}$};
    \foreach \j in {1,...,5} \node[red] at (\j,-7) {$y_{\j}$};
    \node[red] at (5.9,-7) {$\cdots$};
  \end{tikzpicture}
  \caption{Diagonalargument: Die $i$-te Ziffer von $y$ unterscheidet sich von der Diagonalziffer $x_{ii}$.}
  \end{figure}
  Setze nun $y = 0{,}y_1y_2y_3\dots$ mit den Ziffern
  \[
    y_i = \min\bigl(\{1, 2, \dots, 8\} \setminus \{x_{ii}\}\bigr) \in \{1, 2, \dots, 8\}.
  \]
  Da $y$ nur Ziffern aus $\{1, \dots, 8\}$ hat, liegt $y$ in $(0, 1)_\R$, und die Darstellung enthält keine $9$er-Periode. Für jedes $n \in \N$ gilt aber $y_n \neq x_{nn}$, d.\,h.\ $y$ unterscheidet sich von $x_n$ in der $n$-ten Nachkommastelle. Wegen der Eindeutigkeit der Dezimaldarstellung ist also $y \neq x_n$ für alle $n$ -- im Widerspruch dazu, dass die Liste ganz $(0, 1)_\R$ aufzählt.
  \end{proof}
\end{enumerate}
\end{beispiel}

\begin{bemerkung}[Kontinuumshypothese]
Es stellt sich die Frage:
\begin{enumerate}[label=(\arabic*)]
  \item Sei $X \subseteq \R$ mit $|X| = \infty$. Gilt dann bereits $X \approx \N$ oder $X \approx \R$?
\end{enumerate}
Nach Gödel (1939) und Cohen (1963) ist diese Aussage \emph{unabhängig} von den ZF-Axiomen: Man könnte wahlweise (1) oder $\neg$(1) als weiteres Axiom hinzunehmen, ohne einen Widerspruch zu erzeugen.
\end{bemerkung}

% ======================================================================
% Kapitel 2
% ======================================================================
\section*{Kapitel 2: Vektorräume}
\addcontentsline{toc}{part}{Kapitel 2: Vektorräume}

\section{Körper und Vektorräume}

\begin{definition}
Ein \emph{Körper} $K = (K, +, \cdot)$ besteht aus einer Trägermenge $K$ sowie zwei Verknüpfungen, der \emph{Addition} $+$ und der \emph{Multiplikation} $\cdot$,
\[
  +\colon K \times K \to K,\ (a, b) \mapsto a + b, \qquad
  \cdot\colon K \times K \to K,\ (a, b) \mapsto a \cdot b = ab,
\]
welche die folgenden Axiome erfüllen:
\begin{enumerate}[label=(\textup{A\arabic*}),leftmargin=3.2em]
  \item $\forall a, b, c \in K\colon (a + b) + c = a + (b + c)$;
  \item $\forall a, b \in K\colon a + b = b + a$;
  \item Es existiert genau ein Element in $K$, welches mit $0 = 0_K$ bezeichnet wird, für das gilt: $\forall a \in K\colon a + 0 = 0 + a = a$;
  \item Zu jedem $a \in K$ existiert genau ein Element in $K$, welches mit $-a$ bezeichnet wird, für das gilt: $a + (-a) = (-a) + a = 0$.
\end{enumerate}
\begin{enumerate}[label=(\textup{M\arabic*}),leftmargin=3.2em]
  \item $\forall a, b, c \in K\colon a(bc) = (ab)c$;
  \item $\forall a, b \in K\colon ab = ba$;
  \item Es existiert genau ein Element in $K \setminus \{0\}$, welches mit $1 = 1_K$ bezeichnet wird, für das gilt: $\forall a \in K\colon a \cdot 1 = 1 \cdot a = a$;
  \item Zu jedem $a \in K \setminus \{0\}$ existiert genau ein Element in $K$, welches mit $a^{-1} = \frac{1}{a}$ bezeichnet wird, für das gilt: $a \cdot a^{-1} = a^{-1} \cdot a = 1$.
\end{enumerate}
\begin{enumerate}[label=(\textup{D}),leftmargin=3.2em]
  \item $\forall a, b, c \in K\colon (a + b) \cdot c = ac + bc$.
\end{enumerate}
\end{definition}

\begin{bem}
Später: (A1)--(A4) besagen, dass $(K, +)$ eine \emph{abelsche Gruppe} ist; ebenso ist $(K \setminus \{0\}, \cdot)$ eine abelsche Gruppe.
\end{bem}

\begin{lemma}[Folgerungen]
Sei $K$ ein Körper.
\begin{enumerate}[label=(\alph*)]
  \item $\forall a \in K\colon a \cdot 0 = 0 \cdot a = 0$.
  \item $\forall a, b \in K\colon (-a)b = -(ab)$.
  \item $\forall a, b \in K\colon ab = 0 \implies (a = 0 \lor b = 0)$, d.\,h.\ $K$ ist \emph{nullteilerfrei}.
\end{enumerate}
\end{lemma}

\begin{proof}
(a) Mit (A3), (A4), (A1) und (D) gilt
\[
  a \cdot 0 = 0 + a \cdot 0 = \bigl(-(a \cdot 0) + a \cdot 0\bigr) + a \cdot 0 = -(a \cdot 0) + a \cdot (0 + 0) = -(a \cdot 0) + a \cdot 0 = 0.
\]
Wegen (M2) ist auch $0 \cdot a = 0$.

(b) Es gilt $ab + (-a)b = \bigl(a + (-a)\bigr)b = 0 \cdot b = 0$. Aufgrund der Eindeutigkeit des additiven Inversen in (A4) gilt also $(-a)b = -(ab)$.

(c) Sei $ab = 0$ und $a \neq 0$. Dann folgt
\[
  b = 1 \cdot b = (a^{-1}a)b = a^{-1}(ab) = a^{-1} \cdot 0 = 0. \qedhere
\]
\end{proof}

\begin{beispiel}
\leavevmode
\begin{enumerate}[label=(\alph*)]
  \item $\Q$, $\R$ und $\mathbb{C}$ sind Körper.
  \item $\Z$ bildet \emph{keinen} Körper, da außer $1$ und $-1$ kein $a \in \Z \setminus \{0\}$ ein multiplikatives Inverses in $\Z$ besitzt. (Später: $\Z$ ist ein \emph{Ring}.)
  \item $\mathbb{F}_2 = \{0, 1\}$ mit den Verknüpfungstafeln
  \begin{center}
  \begin{tabular}{c|cc}
    $+$ & $0$ & $1$\\\hline
    $0$ & $0$ & $1$\\
    $1$ & $1$ & $0$
  \end{tabular}
  \qquad\qquad
  \begin{tabular}{c|cc}
    $\cdot$ & $0$ & $1$\\\hline
    $0$ & $0$ & $0$\\
    $1$ & $0$ & $1$
  \end{tabular}
  \end{center}
  Mit etwas Geduld überprüft man die Axiome: $\mathbb{F}_2$ ist \emph{der} Körper mit zwei Elementen. Ähnlich (durch Ausprobieren) finden sich Körper $\mathbb{F}_3$, $\mathbb{F}_4$, $\mathbb{F}_5$ mit jeweils $3$, $4$, $5$ Elementen, aber \emph{kein} Körper mit $6$ Elementen.
\end{enumerate}
\end{beispiel}

\begin{definition}
Sei $K$ ein Körper. Für $a \in K$ und eine ganze Zahl $m \in \Z$ definieren wir
\[
  m \cdot a =
  \begin{cases}
    \underbrace{a + a + \dots + a}_{m \text{ Summanden}} & \text{falls } m > 0,\\
    0_K & \text{falls } m = 0,\\
    -\bigl((-m) \cdot a\bigr) & \text{falls } m < 0.
  \end{cases}
\]
Es gelten dann die folgenden Rechenregeln (Übung): Für $m, n \in \Z$ und $a \in K$ ist
\[
  m \cdot a = (m \cdot 1_K) \cdot a, \qquad (m + n) \cdot a = m \cdot a + n \cdot a, \qquad (mn) \cdot a = m \cdot (n \cdot a).
\]
Links steht dabei jeweils die eben definierte Vielfachenbildung, in $(m \cdot 1_K) \cdot a$ dagegen außen die Multiplikation in $K$.
\end{definition}

\begin{definition}
Sei $m \in \N$. Bekanntlich gibt es zu jedem $a \in \Z$ eindeutige $q, r \in \Z$ mit
\[
  a = mq + r \qquad\text{und}\qquad 0 \leq r < m.
\]
Wir bezeichnen den Rest $r$ von $a$ bei Division durch $m$ mit $r_m(a) \in \{0, 1, \dots, m - 1\}$ und setzen
\[
  R_m(\Z) = \operatorname{Bild}(r_m) = \{0, 1, \dots, m - 1\}.
\]
Darauf erklären wir die Verknüpfungen
\begin{align*}
  +_m\colon R_m(\Z) \times R_m(\Z) &\to R_m(\Z), & a +_m b &= r_m(a + b),\\
  \cdot_m\colon R_m(\Z) \times R_m(\Z) &\to R_m(\Z), & a \cdot_m b &= r_m(a \cdot b).
\end{align*}
Hier lassen sich alle Körperaxiome außer (M4) nachprüfen; (M4) gilt jedoch etwa für $m \in \{2, 3, 5, 7\}$.
\end{definition}

\begin{satz}
Sei $m \in \N$. Dann ist $R_m(\Z)$ genau dann ein Körper, wenn $m \in \Pm$ eine Primzahl ist.
\end{satz}

\begin{bem}
Den Beweis liefern wir später (elementare Zahlentheorie).
\end{bem}

\begin{definition}
Sei $K$ ein Körper (beispielsweise $K = \R$ oder $K = \mathbb{F}_2$). Ein \emph{Vektorraum} $V = (V, +, \cdot)$ über $K$, kurz ein \emph{$K$-Vektorraum}, besteht aus einer Trägermenge $V$ sowie zwei Verknüpfungen, einer \emph{Vektoraddition} und einer \emph{skalaren Multiplikation},
\[
  +\colon V \times V \to V,\ (x, y) \mapsto x + y, \qquad
  \cdot\colon K \times V \to V,\ (a, x) \mapsto a \cdot x = ax,
\]
welche den folgenden Bedingungen genügen:
\begin{enumerate}[label=(\textup{VA\arabic*}),leftmargin=3.8em]
  \item $\forall x, y, z \in V\colon x + (y + z) = (x + y) + z$;
  \item $\forall x, y \in V\colon x + y = y + x$;
  \item Es existiert genau ein $+$-neutrales Element, bezeichnet mit $0 = 0_V$, für das also gilt: $\forall x \in V\colon x + 0 = 0 + x = x$;
  \item Zu jedem $x \in V$ existiert genau ein $+$-inverses Element, bezeichnet mit $-x$, für das also gilt: $x + (-x) = 0$.
\end{enumerate}
\begin{enumerate}[label=(\textup{SM\arabic*}),leftmargin=3.8em]
  \item $\forall a, b \in K\ \forall x \in V\colon (ab)x = a(bx)$;
  \item $\forall x \in V\colon 1_K \cdot x = x$;
  \item $\forall a \in K\ \forall x, y \in V\colon a(x + y) = ax + ay$;
  \item $\forall a, b \in K\ \forall x \in V\colon (a + b)x = ax + bx$.
\end{enumerate}
Es gilt \glqq Punkt vor Strich\grqq. Zur Unterscheidung heißen die Elemente von $V$ \emph{Vektoren} und die Elemente von $K$ \emph{Skalare}.
\end{definition}

\begin{lemma}[Folgerungen]
Sei $V$ ein $K$-Vektorraum. Weitere Rechenregeln lassen sich ableiten, zum Beispiel:
\begin{enumerate}[label=(\alph*)]
  \item $\forall a \in K\ \forall x \in V\colon a \cdot 0_V = 0_V = 0_K \cdot x$;
  \item $\forall a \in K\ \forall x \in V\colon (-a)x = -(ax) = a(-x)$, insbesondere $(-1_K)x = -x$;
  \item $\forall a \in K\ \forall x \in V\colon ax = 0_V \implies (a = 0_K \lor x = 0_V)$.
\end{enumerate}
\end{lemma}

\begin{proof}
(a) Nach (VA3) und (SM3) ist $a0_V = a(0_V + 0_V) = a0_V + a0_V$; Addition von $-(a0_V)$ auf beiden Seiten liefert $0_V = a0_V$. Ebenso ist nach (SM4) $0_Kx = (0_K + 0_K)x = 0_Kx + 0_Kx$, also $0_Kx = 0_V$.

(b) Mit (SM4) und (a) gilt $ax + (-a)x = \bigl(a + (-a)\bigr)x = 0_Kx = 0_V$, und mit (SM3) und (a) gilt $ax + a(-x) = a\bigl(x + (-x)\bigr) = a0_V = 0_V$. Wegen der Eindeutigkeit des Inversen in (VA4) folgt $(-a)x = -(ax) = a(-x)$. Für $a = 1_K$ erhalten wir mit (SM2) insbesondere $(-1_K)x = -(1_Kx) = -x$.

(c) Sei $ax = 0_V$ und $a \neq 0_K$. Mit (SM2), (SM1) und (a) folgt
\[
  x = 1_Kx = (a^{-1}a)x = a^{-1}(ax) = a^{-1}0_V = 0_V. \qedhere
\]
\end{proof}

\begin{bem}
Für Körper und für Vektorräume nutzen wir die Kurzschreibweisen $a - b = a + (-b)$ bzw.\ $x - y = x + (-y)$.
\end{bem}

\end{document}
